% \iffalse
%<*driver>
\ProvidesFile{seqenseigne-doc.dtx}%
[2025/10/24 v1.0-dev seqenseigne – Documentation utilisateur]
%</driver>
% \fi
%
% =============================================================
% seqenseigne-doc.dtx
% Module : DOC
% Rôle   :
%   - manuel utilisateur du package
%   - explication des concepts pédagogiques
%   - exemples compilables complets
%   - index des macros et environnements
%
% NOTE :
%   Ici tu colles la grande doc narrative que tu avais commencée :
%   "Le package seqenseigne – structurer des séquences d’enseignement"
%   + les exemples \begin{seqBlocObjectifs} ... \end{seqBlocObjectifs}%
%   + mode d'emploi pour charger les données fusionnées, etc.
%
%   On traite ce fichier comme le .dtx de doc classique LaTeX.
% =============================================================
%
%<*package>
\NeedsTeXFormat{LaTeX2e}%
\ProvidesPackage{seqenseigne-doc}%
[2025/10/24 v1.0-dev seqenseigne documentation module]

% \section{Introduction}%
%  Le paquet \texttt{seqenseigne} est le résultat d'un processus de plusieurs années, lié à ma pratique d'enseignant.
%  En effet, quand j'ai découvert le métier, j'ai été tout de suite confronté à plusieurs problématiques. \par
%  Tout d'abord, comment construire mes séquences? J'ai immédiatement pensé qu'il serait plus efficace de reprendre des
%  contenus déjà existant, éventuellement en les adaptant, plutôt que de tout réinventer\ldots Et de fait, la quantité 
%  de matériel pédagogique disponible est énorme et généralement de très bonne qualité. Je me suis donc centré sur le 
%  processus d'apropriation et d'homogénéisation de ces contenus. \par
%  Ensuite, quel outil logiciel utiliser? La production de documents avec les suites bureautiques dominantes m'a toujours 
%  semblé s'aparenter à du \og bricolage \fg{}. Surtout, ce qui est un réel problème pour les mathématiques, l'inclusion 
%  de formules et de graphiques. Alors j'ai naturellement adopté \LaTeX. \par
%  Enfin, (surtout!) pour quelle utilisation pédagogique? En classe, faire cours aux élèves en leur expliquant quelques exemples 
%  (si j'ai du temps de reste\ldots) et les laisser faire les exercices chez eux, c'est à dire en fait, perpétuer la méthode
%  d'enseignement que j'ai moi-même \og subie \fg{} en tant qu'élève, me semble être une erreur. En effet, je crois que cette méthode convient 
%  essentiellement aux élèves qui ont des facilités de compréhension et tend à exclure les autres au lieu de leur permettre de progresser.
%  J'ai donc choisi d'explorer l'apprentissage inversé, aussi appelé \og classe inversée \fg{}. \bigbreak \par
%  Ainsi, le paquet \texttt{seqenseigne} permet de mettre en forme plusieurs documents : 
%  \begin{seqColItem}[nbCols=1,label=\ding{227}]
%  \item une \textbf{fiche de suivi de séquence} qui permet à l'enseignant de suivre la progression du travail que les élèves font en autonomie 
%    durant les séances et qui récapitule les dates-clés de la séquence. \newline
%    Voir la section \ref{subsec:ficheFic} pour la structure générale du fichier et la section \ref{subsec:ficheOutils} pour les macros et environnements disponibles.
%   \item un \textbf{plan de travail} qui peut contenir les éléments suivants:
%    \begin{seqColItem}%
%     \item la liste des tâches à réaliser en classe
%     \item la liste des tâches à réaliser chez soi
%     \item les dates clés de la séquence
%     \item les ressources disponibles pour réussir
%    \end{seqColItem}%
%    Voir la section \ref{subsec:planFic} pour la structure générale du fichier et la section \ref{subsec:planOutils} pour les macros et environnements disponibles.
%   \item un \textbf{livret de séquence} qui peut contenir les éléments suivants:
%    \begin{seqColItem}[nbCols=1]
%     \item les prérequis pour réussir la séquence
%     \item des exercices de révision
%     \item les objectifs de fin de séquence
%     \item une carte mentale représentant les notions à maîtriser
%     \item les méthodes clés pour résoudre les exercices
%     \item des exercices regroupés en trois séries de difficulté croissante
%     \item les solutions au questionnaire d'autoévaluation initiale et aux exercices
%    \end{seqColItem}%
%    Il n'y a pas de cours (définitions, propriétés, démonstrations, etc) dans le livret de séquence car ces éléments sont préférentiellement produits sous forme de capsules vidéo. \newline
%    Voir la section \ref{subsec:livretFic} pour la structure générale du fichier et la section \ref{subsec:livretOutils} pour les macros et environnements disponibles.
%  \item des \textbf{évaluations bilan} qui peut contenir les éléments suivants:
%    \begin{seqColItem}[nbCols=1]
%     \item des QCM
%     \item des exercices \og classiques \fg{}%
%    \end{seqColItem}%
%    Voir la section \ref{subsec:evalFic} pour la structure générale du fichier et la section \ref{subsec:evalOutils} pour les macros et environnements disponibles.
% \end{seqColItem}%
%
% \newpage
% \section{Comment utiliser \texttt{seqenseigne}}%
%  \subsection{Installation de \texttt{seqenseigne}}%
%   Le fichier \texttt{seqenseigne.sty} doit être placé dans un répertoire visible pour \LaTeX:
%   \begin{seqColItem}[nbCols=1]
%    \item dans le même répertoire que le document
%    \item dans l'arborescence \LaTeX\, de l'utilisateur. \newline Par exemple, le répertoire \texttt{\textasciitilde/.textive/texmf-var/arch/tlpkg} pour TexLive.
%    \item dans l'arborescence \LaTeX\, système. \newline Par exemple, le répertoire \texttt{/usr/local/share/texmf/tex/latex} pour TexLive.
%   \end{seqColItem}%
%
%  \subsection{Structure des documents utilisant \texttt{seqenseigne}}%
% 	\subsubsection{Fiche de suivi} \label{subsec:ficheFic}%
%    La fiche de suivi comprend une ou deux pages, en fonction du nombre d'élèves de la classe. La structure du fichier \LaTeX\, est très simple:\newline
%    \verb|  \documentclass[a4paper,11pt]{article}|\newline
%    \verb|  \usepackage{seqenseigne}|\newline
%    \verb|  \begin{document}|\newline
%    \verb|   \begin{titreFiche}[numSeq=1,titreSeq=Théorème de Pythagore,| \newline
%    \verb|                      etab=Collège Exemple,classe=4\ieme{}A]|\newline
%    \verb|    \seq@ligne@eleve{Laurent COSTE}|\newline
%    \verb|   \end{titreFiche}|\newline
%    \verb|  \end{document}|\newline
%    Les différents éléments de la fiche de suivi sont détaillés à partir de la page \pageref{elementsSuivi}.
%
% 	\subsubsection{Plan de travail} \label{subsec:planFic}%
%    Le plan de travail comprend une seule page. La structure du fichier \LaTeX\, est très simple:\newline
%    \verb|  \documentclass[a4paper,11pt]{article}|\newline
%    \verb|  \usepackage{seqenseigne}|\newline
%    \verb|  \begin{document}|\newline
%    \verb|   \begin{titrePlan}[numSeq=1,titreSeq=Théorème de Pythagore,| \newline
%    \verb|                      etab=Collège Exemple,classe=4\ieme{}A]|\newline
%    \verb|    \begin{seqBloc}{En classe}| \newline
%    \verb|     \item Je travaille!| \newline
%    \verb|    \end{seqBloc}| \newline
%    \verb|    &| \newline
%    \verb|    \begin{seqBloc}{Chez moi}| \newline
%    \verb|    \item Je travaille| \newline
%    \verb|     \begin{blocN2} | \newline
%    \verb|      \item Le soir| \newline
%    \verb|      \item Le week-end| \newline
%    \verb|     \end{blocN2}| \newline
%    \verb|    \item Je m'amuse aussi!| \newline
%    \verb|   \end{seqBloc}| \newline
%    \verb|   \newline| \newline
%    \verb|  \end{titrePlan}| \newline
%    \verb|  \end{document}|\newline
%    Les différents éléments du plan de travail sont détaillés à partir de la page \pageref{elementsPlan}.
%    
% 	\subsubsection{livret de séquence} \label{subsec:livretFic}%
%    Plusieurs livrets de séquence peuvent coexister dans le même répertoire. \newline
%    Je conseille fortement une structure modulaire comprenant un document-maître et des sous-documents, telle que:\newline
%    \verb|  \documentclass[a4paper,11pt]{article}|\newline
%    \verb|  \usepackage{seqenseigne}|\newline
%    \verb|  \begin{document}|\newline
%    \verb|   \titreLivret[theme=geometrie,titreSeq=Théorème de Pythagore,|\newline
%    \verb|                      niveau=quatrième,numseq=1}|\newline
%    \verb|    \input{Include/Prerequis.tex}|\newline
%    \verb|    \input{./include/Seq01AutoEvaluation.tex}|\newline
%    \verb|   \seqTitresection{Objectifs et connaissances}|\newline
%    \verb|    \input{Include/Objectifs.tex}|\newline
%    \verb|    \input{Include/CarteMentale.tex}|\newline
%    \verb|   \seqTitresection{Méthodes}|\newline
%    \verb|    \input{Include/Methodes.tex}|\newline
%    \verb|   \seqTitresection{Exercices}|\newline
%    \verb|    \input{Include/SerieFondamentale.tex}|\newline
%    \verb|    \input{Include/SerieAvancee.tex}|\newline
%    \verb|    \input{Include/SerieExploration.tex}|\newline
%    \verb|   \seqAfficheCorriges|\newline
%    \verb|  \end{document}|\newline
%    Dans l'exemple précédent, il faut créer manuellement le sous-répertoire \meta{nom du fichier}\texttt{Include} qui contiendra tous les sous-documents.\newline
%    Par ailleurs, je conseille aussi fortement de placer les images incluses dans les différentes sections du livret de séquence dans un répertoire 
%    \meta{nom du fichier}\texttt{Images}, à créer manuellement. \newline
%    Enfin, la compilation génère un fichier de corrigé par exercice. Tous ces fichiers sont stockés dans un sous-répertoire \meta{nom du fichier}\texttt{Corriges} qu'il faut créer manuellement.\newline
%    Les différents éléments du livret de séquence sont détaillés à partir de la page \pageref{elementslivret}.
%
% 	\subsubsection{\'Evaluation bilan} \label{subsec:evalFic}%
%    Plusieurs évaluations peuvent coexister dans le même répertoire. \newline
%    Je conseille fortement une structure modulaire comprenant un document-maître et des sous-documents, telle que:\newline
%    \verb|  \documentclass[a4paper,11pt]{article}|\newline
%    \verb|  \usepackage{seqenseigne}|\newline
%    \verb|  \usepackage{pgfplots}|\newline
%    \verb|  \usepackage{eurosym}|\newline
%    \verb|  \begin{document}|\newline
%    \verb|  \titreEval[numSeq=03,titreSeq=Proportionnalité,niveauDiff=fondamental,|\newline
%    \verb|             etab=Collège exemple,classe=4\ieme{}A]|\newline
%    \verb|  \begin{bareme}|\newline
%    \verb|  	\baremeItem{QCM}{5 points}|\newline
%    \verb|  	\baremeItem{Exercice II}{6 points}|\newline
%    \verb|  	\baremeItem{Présentation de la copie et usage de la langue française}{1 point}|\newline
%    \verb|  \end{bareme}|\newline
%    \verb|  \input{Include/exoI.tex}|\newline
%    \verb|  \input{Include/exoII.tex}|\newline
%    \verb|  \afficheCorrigesEval|\newline
%    \verb|  \end{document}|\newline
%    Dans l'exemple précédent, il faut créer manuellement le sous-répertoire \meta{nom du fichier}\texttt{Include} qui contiendra tous les sous-documents.\newline
%    Par ailleurs, je conseille aussi fortement de placer les images incluses dans les différentes sections du livret de séquence dans un répertoire 
%    \meta{nom du fichier}\texttt{Images}, à créer manuellement. \newline
%    Enfin, la compilation génère un fichier de corrigé par exercice. Ils sont stockés dans un sous-répertoire \meta{nom du fichier}\texttt{Corriges} qu'il faut créer manuellement.\newline
%    Les différents éléments de l'évaluation bilan sont détaillés à partir de la page \pageref{elementsEval}.
%
%  \subsection{Paquets chargés par \texttt{seqenseigne}}%
%   Le paquet \texttt{seqenseigne} charge les paquets suivants:
%   \begin{seqColItem}%
%    \item \texttt{inputenc} avec l'option \texttt{utf8}%
%    \item \texttt{fontenc} avec l'option \textsf{T1}%
%    \item \texttt{lmodern}%
%    \item \texttt{babel} avec l'option \texttt{french}%
%    \item \texttt{geometry}%
%    \item \texttt{fancyhdr}%
%    \item \texttt{lastpage}%
%    \item \texttt{multicol}, \texttt{multirow}%
%    \item \texttt{amsmath}, \texttt{amsfonts}, \texttt{amssymb}%
%    \item \texttt{tcolorbox} avec les options \texttt{theorems}, \texttt{breakable} et \texttt{skins}%
%    \item \texttt{tabularray} avec les bibliothèques \texttt{booktabs} et \texttt{longtable}%
%    \item \texttt{tabularx}, \texttt{xltabular}, \texttt{booktabs}%
%    \item \texttt{xcolor}, \texttt{colortbl}%
%    \item \texttt{xkeyval}, \texttt{xifthen}, \texttt{etoolbox}, \texttt{kvoptions}%
%    \item \texttt{forloop}%
%    \item \texttt{datatool} (chargement des bases de données \texttt{.dbtex})
%    \item \texttt{tikz} avec ses librairies \texttt{babel}, \texttt{intersections} et \texttt{mindmap}%
%    \item \texttt{tkz-tab}, \texttt{tkz-base}, \texttt{tkz-euclide}%
%    \item \texttt{hyperref}%
%    \item \texttt{scratch3} (blocs Scratch dans les exercices)
%    \item \texttt{eurosym}, \texttt{pifont}%
%    \item \texttt{enumitem} avec l'option \textsf{inline}%
%    \item \texttt{datetime}, \texttt{environ}, \texttt{xintexpr}%
%   \end{seqColItem}%
%   La documentation de ces paquets est disponible sur \texttt{CTAN}. \newline
%   Il est possible d'appeler tout autre paquet nécessaire à la mise en forme du document en utilisant la macro |\usepackage|.
%
%  \subsection{\`A propos de \TikZ}%
%   J'ai choisi d'utiliser \TikZ pour tous les graphiques. Sa librairie \texttt{mindmap} permet de construire les cartes mentales. \newline
%   Le paquet \texttt{tkz-euclide} sert pour les figures en géométrie euclidienne.\newline
%   Le paquet \texttt{tkz-tab} sert pour les tableaux de variation.
%   \medskip \newline
%   Voici quelques macro pour faciliter le tracé des graphiques.
%   \medskip \newline
%   \DescribeMacro{\seqOrigine}%
%    La macro |\seqOrigine| place le point $O$ à l'origine du repère.
%    \medskip \newline
%   \DescribeMacro{\seqAxeX}%
%    La macro |\seqAxeX|\oarg{options} trace l'axe des abscisses.\newline
%    L'option \texttt{ord}=\meta{nombre} est l'ordonnée de l'axe. Elle vaut 0 par défaut. \newline
%    L'option \texttt{absMin}=\meta{nombre} est l'abscisse minimum. Cette option est obligatoire. \newline
%    L'option \texttt{absMax}=\meta{nombre} est l'abscisse maximum. Cette option est obligatoire. \newline
%    L'option \texttt{abslist}=\{\meta{liste}\} est la liste des abscisses qui seront marquées sur l'axe. Elle est vide par défaut. \newline
%    L'option \texttt{libelle}=\meta{texte} est le libellé placé au dessus à droite de l'axe. Il est vide par défaut. 
%    \medskip \newline
%   \DescribeMacro{\seqAxeY}%
%    La macro |\seqAxeY|\oarg{options} trace l'axe des ordonnées.\newline
%    L'option \texttt{abs}=\meta{nombre} est l'abscisse de l'axe. Elle vaut 0 par défaut. \newline
%    L'option \texttt{ordMin}=\meta{nombre} est l'ordonnée minimum. Cette option est obligatoire. \newline
%    L'option \texttt{ordMax}=\meta{nombre} est l'ordonnée maximum. Cette option est obligatoire. \newline
%    L'option \texttt{ordlist}=\{\meta{liste}\} est la liste des ordonnées qui seront marquées sur l'axe. Elle est vide par défaut. \newline
%    L'option \texttt{libelle}=\meta{texte} est le libellé placé au dessus à droite de l'axe. Il est vide par défaut. 
%    \medskip \newline
%   Par exemple: \newline
%   \verb|\begin{tikzpicture}[scale=0.5]| \newline
%   \verb| \seqAxeX[absMin=-4,absMax=5,absList={-2,2,4},libelle=nombre]| \newline
%	\verb| \seqAxeY[ordMin=-1,ordMax=5,ordList={1,2,3},libelle=prix]| \newline
%	\verb| \seqOrigine| \newline
%	\verb|\end{tikzpicture}| \newline
%   donne:
%   \begin{center} \begin{tikzpicture}[scale=0.5]
%    \seqAxeX[absMin=-4,absMax=5,absList={-2,2,4},libelle=nombre] \seqAxeY[ordMin=-1,ordMax=5,ordList={1,2,3},libelle=prix] \seqOrigine
%   \end{tikzpicture} \end{center}%
%
%  \subsection{\`A propos de \textsf{tcolorbox}}%
%   Le paquet \textsf{tcolorbox} permet de définir des boîtes. Je l'ai utilisé pour délimiter les différentes parties des documents.\newline
%   Les boîtes de titre de chaque document (fiche de suivi, plan de travail et livret de séquence) ont une hauteur fixe; les hauteurs des autres boîtes s'adaptent à leur contenu.\newline
%   Les boîtes des numéros d'exercices ont une largeur fixe; les largeurs des autres boîtes s'adaptent à leur contenu.\newline 
%   Les boîtes qui contiennent chaque série d'exercice peuvent s'étendre sur plusieurs pages; les autres sont d'un seul bloc, c'est à dire sur une seule page.\newline
%   Attention aux débordements si elles sont trop grandes!
%
%   \subsection{Définition des couleurs} \label{subsec:themes}%
%    Pour le thème \og \textit{Nombres et calcul} \fg{}, les couleurs de la maquette sont:
%     \begin{tcolorbox}[colback=\seqCouleurRougeVif] |\seqCouleurRougeVif| défini par \textit{red!85!black}%
%     \end{tcolorbox}%
%     \begin{tcolorbox}[colback=\seqCouleurRougeModere] |\seqCouleurRougeModere| défini par \textit{red!50!white}%
%     \end{tcolorbox}%
%     \begin{tcolorbox}[colback=\seqCouleurRougeMoyen] |\seqCouleurRougeMoyen| défini par \textit{red!25!white}%
%     \end{tcolorbox}%
%     \begin{tcolorbox}[colback=\seqCouleurRougePale] |\seqCouleurRougePale| défini par \textit{red!10!white}%
%     \end{tcolorbox}%
%     \begin{tcolorbox}[colback=\seqCouleurJauneModere] |\seqCouleurJauneModere| défini par \textit{yellow!50!white}%
%     \end{tcolorbox}%
%     \begin{tcolorbox}[colback=\seqCouleurJaunePale] |\seqCouleurJaunePale| défini par \textit{yellow!10!white}%
%     \end{tcolorbox}%
%    Pour le thème \og \textit{Organisation des données} \fg{}, les couleurs de la maquette sont:
%     \begin{tcolorbox}[colback=\seqCouleurBleuVif] |\seqCouleurBleuVif| défini par \textit{blue!85!black}%
%     \end{tcolorbox}%
%     \begin{tcolorbox}[colback=\seqCouleurBleuModere] |\seqCouleurBleuModere| défini par \textit{blue!50!white}%
%     \end{tcolorbox}%
%     \begin{tcolorbox}[colback=\seqCouleurBleuMoyen] |\seqCouleurBleuMoyen| défini par \textit{blue!25!white}%
%     \end{tcolorbox}%
%     \begin{tcolorbox}[colback=\seqCouleurBleuPale] |\seqCouleurBleuPale| défini par \textit{blue!10!white}%
%     \end{tcolorbox}%
%     \begin{tcolorbox}[colback=\seqCouleurOrangeModere] |\seqCouleurOrangeModere| défini par \textit{orange!50!white}%
%     \end{tcolorbox}%
%     \begin{tcolorbox}[colback=\seqCouleurOrangePale] |\seqCouleurOrangePale| défini par \textit{orange!10!white}%
%     \end{tcolorbox}%
%    Pour le thème \og \textit{Espace et géométrie} \fg{}, les couleurs de la maquette sont:
%     \begin{tcolorbox}[colback=\seqCouleurVioletVif] |\seqCouleurVioletVif| défini par \textit{violet!85!black}%
%     \end{tcolorbox}%
%     \begin{tcolorbox}[colback=\seqCouleurVioletModere] |\seqCouleurVioletModere| défini par \textit{violet!50!white}%
%     \end{tcolorbox}%
%     \begin{tcolorbox}[colback=\seqCouleurVioletMoyen] |\seqCouleurVioletMoyen| défini par \textit{violet!25!white}%
%     \end{tcolorbox}%
%     \begin{tcolorbox}[colback=\seqCouleurVioletPale] |\seqCouleurVioletPale| défini par \textit{violet!10!white}%
%     \end{tcolorbox}%
%     \begin{tcolorbox}[colback=\seqCouleurVertModere] |\seqCouleurVertModere| défini par \textit{green!50!white}%
%     \end{tcolorbox}%
%     \begin{tcolorbox}[colback=\seqCouleurVertPale] |\seqCouleurVertPale| défini par \textit{green!10!white}%
%     \end{tcolorbox}%
%    Quand il n'y a pas de thème, les couleurs de la maquette sont:
%     \begin{tcolorbox}[colback=\seqCouleurNoirVif,coltext=white] |\seqCouleurNoirVif| défini par \textit{black!85!black}%
%     \end{tcolorbox}%
%     \begin{tcolorbox}[colback=\seqCouleurNoirModere,coltext=white] |\seqCouleurNoirModere| défini par \textit{black!50!white}%
%     \end{tcolorbox}%
%     \begin{tcolorbox}[colback=\seqCouleurNoirMoyen] |\seqCouleurNoirMoyen| défini par \textit{black!25!white}%
%     \end{tcolorbox}%
%     \begin{tcolorbox}[colback=\seqCouleurNoirPale] |\seqCouleurNoirPale| défini par \textit{black!10!white}%
%     \end{tcolorbox}%
%     \begin{tcolorbox}[colback=\seqCouleurGrisModere] |\seqCouleurGrisModere| défini par \textit{gray!50!white}%
%     \end{tcolorbox}%
%     \begin{tcolorbox}[colback=\seqCouleurGrisPale] |\seqCouleurGrisPale| défini par \textit{gray!10!white}%
%     \end{tcolorbox}%

%  \subsection{Environnements de listes}%
%   \DescribeEnv{colEnum}%
%    L'environnement |colEnum|\oarg{options} est destiné à combiner |\enumerate| et |\multicols|.\newline
%    L'option \texttt{nbCols}=\meta{nombre} sert à fixer le nombre de colonnes. Elle vaut 2 par défaut.
%   \medskip \newline
%   \DescribeEnv{colItem}%
%    L'environnement |colItem|\oarg{options} est destiné à combiner |\itemize| et |\multicols|.\newline
%    L'option \texttt{nbCols}=\meta{nombre} sert à fixer le nombre de colonnes. Elle vaut 2 par défaut.\newline
%    L'option \texttt{label}=\meta{texte} sert à fixer la puce devant chaque item. Elle vaut \verb|$\bullet$| par défaut.
%   \medskip \newline
%   Exemple:
%    \begin{seqColEnum}%
%     \item Enum 1
%     \item Enum 2
%      \begin{seqColItem}[nbCols=1]
%       \item Item 2.1
%       \item Item 2.2
%      \end{seqColItem}%
%     \item Enum 3
%     \item Enum 4
%    \end{seqColEnum}%
%   On peut imbriquer ces deux environnements.
%
% \newpage
% \section{Les différents éléments de la fiche de suivi de séquence} \label{subsec:ficheOutils} \label{elementsSuivi}%
%  \subsection{Titre de la fiche de suivi}%
%   \DescribeEnv{titreFiche}%
%    L'environnement |titreFiche|\oarg{options} met en forme le titre du document.\newline
%    Il débute également un tableau à dix colonnes dans lequel seront listés les élèves à l'aide de la macro |\seq@ligne@eleve|. \newline
%    Il définit enfin les entêtes et pieds de page.\newline
%    L'option obligatoire \texttt{numSeq}=\meta{texte} est le numéro de la séquence dans la progression. \newline
%    L'option obligatoire \texttt{titreSeq}=\meta{texte} est le nom de la séquence. \newline
%    L'option obligatoire \texttt{etab}=\meta{texte} est le nom de l'établissement. \newline
%    L'option obligatoire \texttt{classe}=\meta{texte} est le nom de la classe. \newline
%    L'option facultative \texttt{année}=\meta{texte} est l'année scolaire. Elle est vide par défaut. \newline
%    L'option facultative \texttt{enseignant}=\meta{texte} est le nom de l'enseignant. Elle est vide par défaut. \newline
%	 Il n'y a pas d'option pour le theme de couleurs qui est obligatoirement gris. \medskip \newline
%    
%   \DescribeMacro{\seq@ligne@eleve}%
%    La macro |\seq@ligne@eleve|\marg{eleve} met en forme une ligne du tableau contenant les élèves de la classe. \newline
%    L'argument \meta{eleve} est obligatoire.
%    \medskip \newline
%    Par exemple, \newline
%    \verb|  	\begin{titreFiche}[numSeq=01,titreSeq=Théorème de Pythagore,|\newline
%    \verb|  	                   etab=Collège Exemple,classe=4\ieme{}A,|\newline
%    \verb|  	                   annee=2017/2018,enseignant=L. Coste]|\newline
%    \verb|  		\seq@ligne@eleve{Laurent COSTE}|\newline
%    \verb|  	\end{titreFiche}|\newline
%    donne :
%    \begin{titreFiche}[numSeq=01,titreSeq=Théorème de Pythagore,etab=Collège Exemple,classe=4\ieme{}A,
%                       annee=2017/2018,enseignant=L. Coste]
%     \seq@ligne@eleve{Laurent COSTE}%
%    \end{titreFiche}%
%
% \newpage
% \section{Les différents éléments du plan de travail} \label{subsec:planOutils} \label{elementsPlan}%
%  \subsection{Titre du plan de travail}%
%   \DescribeEnv{titrePlan}%
%    L'environnement |titrePlan|\oarg{options} met en forme le titre du document.\newline
%    Il débute également un tableau à deux colonnes dans lequel viendront se placer les différents blocs.
%    Il définit enfin les entêtes et pieds de page.\newline
%    L'option obligatoire \texttt{numSeq}=\meta{texte} est le numéro de la séquence dans la progression. \newline
%    L'option obligatoire \texttt{titreSeq}=\meta{texte} est le nom de la séquence. \newline
%    L'option obligatoire \texttt{etab}=\meta{texte} est le nom de l'établissement. \newline
%    L'option obligatoire \texttt{classe}=\meta{texte} est le nom de la classe. \newline
%    L'option facultative \texttt{année}=\meta{texte} est l'année scolaire. Elle est vide par défaut. \newline
%    L'option facultative \texttt{enseignant}=\meta{texte} est le nom de l'enseignant. Elle est vide par défaut. \newline
%    L'option facultative \texttt{theme}=\meta{texte} définit le thème de couleurs utilisé.
%     Les différentes valeurs possibles sont \textit{nombres}, \textit{donnees}, \textit{geometrie} et \textit{gris} (valeur par défaut). 
%     Voir \ref{subsec:themes} pour leur définition.
%
%  \subsection{Blocs}%
%   \DescribeEnv{bloc}%
%    L'environnement |bloc|\marg{titre} construit une boîte comportant un titre, 
%    dans laquelle seront listés des éléments en utilisant la macro |\item|.\newline
%    L'argument \meta{titre} est obligatoire.
%    \medskip \newline
%   \DescribeEnv{blocN2}%
%    L'environnement |blocN2| permet d'imbriquer un deuxième niveau dans le bloc en utilisant également la macro |\item|.
%    \medskip \newline
%    Par exemple,\newline
%    \verb|  \begin{titrePlan}[numSeq=03,titreSeq=Proportionnalité,annee=2017/2018,| \newline
%    \verb|      etab=Collège exemple,classe=4\ieme{}A,enseignant=L. Coste,theme=nombres]| \newline
%    \verb|   \begin{seqBloc}{En classe}| \newline
%    \verb|    \item Je fais les exercices de révivion| \newline
%    \verb|    \item Je m'exerce en faisant les exercices des trois séries| \newline
%    \verb|    \item Je pose les questions sur les difficultés rencontrées lors de mon | \newline
%    \verb|      travail en classe ou chez moi| \newline
%    \verb|   \end{seqBloc}| \newline
%    \verb|   &| \newline
%    \verb|   \begin{seqBloc}{Chez moi}| \newline
%    \verb|    \item J'apprends mon cours et je fais ma fiche de résumé| \newline
%    \verb|     \begin{blocN2} | \newline
%    \verb|      \item Cours et savoir-faire pages 122 à 125 | \newline
%    \verb|      \item Carte mentale page 128| \newline
%    \verb|     \end{blocN2}| \newline
%    \verb|    \item Je fais les exercices corrigés du manuel| \newline
%    \verb|     \begin{blocN2}| \newline
%    \verb|      \item \og Faire le point \fg{} page 128| \newline
%    \verb|      \item Exercices numérotés sur fond vert| \newline
%    \verb|     \end{blocN2}| \newline
%    \verb|   \end{seqBloc}| \newline
%    \verb|   \newline| \newline
%    \verb|  \end{titrePlan}| \newline
%    donne : 
%    \newpage
%    \begin{titrePlan}[numSeq=03,titreSeq=Proportionnalité,etab=Collège exemple,classe=4\ieme{}A,
%                      annee=2017/2018,enseignant=L. Coste,theme=nombres]
%    	\begin{seqBloc}{En classe}%
%    		\item Je fais les exercices de révision
%    		\item Je m'exerce en faisant les exercices des trois séries
%    		\item Je pose les questions sur les difficultés rencontrées lors de mon travail en classe ou chez moi
%    	\end{seqBloc}%
%    	&
%    	\begin{seqBloc}{Chez moi}%
%    		\item J'apprends mon cours et je fais ma fiche de résumé
%    			\begin{blocN2} 
%    				\item Cours et savoir-faire pages 122 à 125 
%    				\item Carte mentale page 128
%    			\end{blocN2}%
%    		\item Je fais les exercices corrigés du manuel
%    			\begin{blocN2}%
%    				\item \og Faire le point \fg{} page 128
%    				\item Exercices numérotés sur fond vert
%    			\end{blocN2}%
%    	\end{seqBloc}%
%    	\\
%    \end{titrePlan}%
%
% \newpage
% \section{Les différents éléments du livret de séquence} \label{subsec:livretOutils} \label{elementslivret}%
%  Si la maquette du livret de séquence est libre, dans le sens où aucun élément n'est obligatoire et l'ordre des éléments présent est indifférent, 
%  les macros et les environnements ont été pensés pour être ordonnés d'une certaine manière.\newline
%  Certains environnements intègrent un titre à la boîte qu'ils construisent. Ils sont destinés à être utilisés une seule fois
%  dans le document. Par exemple: la carte mentale. Ils n'est pas nécessaire d'insérer un titre de section avant ces environnements.\newline
%  A contrario, d'autres environnements n'intrègent pas de titre: ils nécessitent un titre de section. Ils sont destinés à être
%  utilisés plusieurs fois à la suite dans le document. Par exemple: les méthodes clés.
%
%  \subsection{Titre du livret de séquence}%
%   \DescribeMacro{\titreLivret}%
%    La macro |\titreLivret|\oarg{options} met en forme le titre du livret de séquence.\newline
%    Elle définit les entêtes et pieds de page.\newline
%    Elle crée les compteurs pour numéroter automatiquement les exercices de révision, les méthodes, les exercices.\newline
%    L'option \texttt{numSeq}=\meta{texte} est le numéro de la séquence dans la progression. \newline
%    L'option \texttt{titreSeq}=\meta{texte} est le nom de la séquence. \newline
%    L'option \texttt{niveau}=\meta{texte} est le niveau de classe de cette séquence. \newline
%    Ces trois options sont obligatoires. \newline
%    L'option facultative \texttt{theme}=\meta{texte} définit le thème de couleurs utilisé.
%     Les différentes valeurs possibles sont \textit{nombres}, \textit{donnees}, \textit{geometrie} et \textit{gris} (valeur par défaut). 
%     Voir \ref{subsec:themes} pour leur définition.
%   \medskip \newline
%   Par exemple : \newline
%   \verb|\titreLivret[numSeq=01,titreSeq=Théorème de Pythagore,niveau=quatrième,|\newline
%   \verb|                theme=geometrie]|\newline
%   donne :
%	\par\vspace{20pt}%
%	\titreLivret[numSeq=01,titreSeq=Théorème de Pythagore,niveau=quatrième,theme=geometrie]
%
%  \subsection{Titres de section}%
%   \DescribeMacro{\seqTitresection}%
%    La macro |\seqTitresection|\marg{titre} sert à inserer un titre au début d'une nouvelle section. \newline
%    Il est placé au début d'une nouvelle page. \newline
%    L'argument \meta{titre} est obligatoire.
%    \medskip \newline
%   Par exemple :\newline
%   \verb|\seqTitresection{Exercices}| \newline
%   donne : 
%	\par\vspace{20pt}%
%   \seqTitresection{Exercices}%
%
%  \subsection{Prérequis et objectifs}%
%   \DescribeEnv{prerequis}%
%    L'environnement |prerequis|\oarg{description} construit une boîte comportant un titre et une courte description, 
%    dans laquelle seront listés les prérequis de la séquence en utilisant la macro |\item|.\newline
%    L'argument optionnel \meta{description} permet de modifier la description par défaut.
%   \medskip \newline
%   \DescribeEnv{prerequisN2}%
%    L'environnement |prerequisN2| permet d'imbriquer un deuxième niveau de prérequis en utilisant également la macro |\item|.
%   \medskip \newline
%   Par exemple:\newline
%   \verb|  \begin{seqPrerequis}|\newline
%   \verb|   \item Carré et racine carrée d'un nombre|\newline
%   \verb|   \item Notions fondamentales concernant les triangles|\newline
%   \verb|  \end{seqPrerequis}|\newline
%   donne :
%   \begin{seqPrerequis}%
%    \item Carré et racine carrée d'un nombre
%    \item Notions fondamentales concernant les triangles
%   \end{seqPrerequis}%
%
%   \DescribeEnv{objectifs}%
%    L'environnement |objectifs|\oarg{description} fonctionne comme l'environnement \newline
%    |prerequis|\oarg{description} ci-avant.
%   \medskip \newline
%   \DescribeEnv{objectifsN2}%
%    L'environnement |objectifsN2| fonctionne comme l'environnement |prerequisN2| ci-avant.
%   \medskip \newline
%    Par exemple:\newline
%    \verb|  \begin{seqObjectif}|\newline
%    \verb|   \item Caractériser le triangle rectangle par l'égalité de Pythagore|\newline
%    \verb|   \item Calculer la longueur d'un côté d'un triangle rectangle|\newline
%    \verb|   \item Démontrer qu'un triangle est rectangle ou n'est pas rectangle|\newline
%    \verb|         en connaissant les longueurs de ses côtés|\newline
%    \verb|  \end{seqObjectif}|\newline
%    donne :
%    \begin{seqObjectif}%
%     \item Caractériser le triangle rectangle par l'égalité de Pythagore
%     \item Calculer la longueur d'un côté d'un triangle rectangle
%     \item Démontrer qu'un triangle est rectangle ou n'est pas rectangle en connaissant les longueurs de ses côtés
%    \end{seqObjectif}%
%
%  \subsection{Révisions}%
%   \DescribeEnv{autoEval}%
%    L'environnement |autoEval|\oarg{description} construit une boîte comportant un titre et une courte description,
%    dans laquelle seront regroupées les exercices de révision.\newline
%    L'argument optionnel \meta{description} permet de modifier la description par défaut.
%   \medskip \newline
%   \DescribeEnv{exoAE}%
%    L'environnement |exoAE| permet de mettre en forme un exercice révision. La numérotation des exercices est automatique.
%   \medskip \newline
%   \DescribeMacro{\corrigeAE}%
%    La macro |\corrigeAE|\marg{corrige} permet de donner le corrigé de l'exercice.
%    Elle se place à l'intérieur de l'environnement |\exoAE|.\newline
%    L'argument obligatoire \meta{corrige} est le corrigé de l'exercice.\newline
%    Tous les corrigés sont automatiquement regroupés en fin de document.
%   \medskip \newline
%    Par exemple:\newline
%    \verb|  \begin{autoEval}|\newline
%    \verb|   \begin{exoAE}|\newline
%    \verb|    Calculer le carré de 16.|\newline
%    \verb|    \corrigeAE{ Le carré de 16 se note $16^2$. Il vaut $16 \times 16 = 256$. }|\newline
%    \verb|   \end{exoAE}|\newline
%    \verb|   \begin{exoAE}|\newline
%    \verb|    Donner le nom de l'hypoténuse du rectangle $RST$.|\newline
%    \verb|    \begin{center} \begin{tikzpicture}[scale=0.5]|\newline
%    \verb|     \coordinate (R) at (0,0); \coordinate (S) at (3,0);|\newline
%    \verb|     \coordinate (T) at (0,5);|\newline
%    \verb|     \draw (0,0) rectangle (0.2,0.2);|\newline
%    \verb|     \draw (R) node[below left]{$R$} -- (S) node[below right]{$S$} --|\newline
%    \verb|           (T) node[above]{$T$} -- cycle;|\newline
%    \verb|    \end{tikzpicture} \end{center}|\newline
%    \verb|    \corrigeAE{ L'hypoténuse de $RST$ est le côté à l'opposé de l'angle droit.|\newline
%    \verb|        C'est aussi le côté le plus long du triangle. Ici, c'est $[ST]$. }|\newline
%    \verb|   \end{exoAE}|\newline
%    \verb|  \end{autoEval}|\newline
%    donne :
%    \begin{autoEval}%
%     \begin{exoAE}%
%      Calculer le carré de 16.
%      \corrigeAE{Le carré de 16 se note $16^2$. Il vaut $16 \times 16 = 256$.}%
%     \end{exoAE}%
%     \begin{exoAE}%
%      Donner le nom de l'hypoténuse du rectangle $RST$.
%      \begin{center} \begin{tikzpicture}[scale=0.5]
%         \coordinate (R) at (0,0);	\coordinate (S) at (3,0);	\coordinate (T) at (0,5);
%         \draw (0,0) rectangle (0.2,0.2);
%         \draw (R) node[below left]{$R$} -- (S) node[below right]{$S$} -- (T) node[above]{$T$} -- cycle;
%      \end{tikzpicture} \end{center}%
%      \corrigeAE{%
%       L'hypoténuse de $RST$ est le côté à l'opposé de l'angle droit. C'est aussi le côté le plus long du triangle. 
%       Ici, c'est $[ST]$.
%      }%
%     \end{exoAE}%
%    \end{autoEval}%
%
%  \subsection{Carte mentale}%
%   \DescribeEnv{carteMentale}%
%    L'environnement |carteMentale|\oarg{description} construit une boîte comportant un titre et une courte description,
%    dans laquelle sera placée la carte mentale représentant de manière synthétique les connaissances de la séquence.\newline
%    L'argument optionnel \meta{description} permet de modifier la description par défaut.\newline
%    La carte mentale elle-même est mise en forme avec l'extension \textsf{mindmap} de \TikZ.
%    \medskip \newline
%   Par exemple:\newline
%    \verb|\begin{seqCarteMentale}|\newline
%    \verb| \begin{center} \begin{tikzpicture}[small mindmap,concept color=\seqCouleurVioletVif,|\newline
%    \verb|    level 1 concept/.style={concept color=\seqCouleurVioletModere},|\newline
%    \verb|    every annotation/.style={fill=\seqCouleurVioletPale}]|\newline
%    \verb|  \node [concept] {Théorème de Pythagore} [counterclockwise from=0]|\newline
%    \verb|   child {node[concept] (egalite) {\'Egalité de Pythagore} }|\newline
%    \verb|   child {node[concept] (hypotenuse) {Hypoténuse d'un triangle rectangle} }|\newline
%    \verb|   child {node[concept] (racine) {Racine carrée} }|\newline
%    \verb|   child {node[concept] (theoreme) {Le théorème} }|\newline
%    \verb|   child {node[concept] (reciproque) {La réciproque du théorème} };|\newline
%    \verb|  \node [annotation,above left=5pt] at (racine.north west){|\newline
%    \verb|    La racine carrée du nombre positif $a$ est le nombre positif|\newline
%    \verb|    $b$ tel que $b^2=a$.};|\newline
%    \verb|  \node [annotation,left=5pt] at (racine.west){|\newline
%    \verb|    Le symbole de la  racine carrée est $\sqrt{\phantom{ab}}$.};|\newline
%    \verb|  \node [annotation,above left=5pt] at (hypotenuse.north){|\newline
%    \verb|    \begin{tikzpicture}[scale=0.5, sharp corners]|\newline
%    \verb|      \coordinate (A) at (0,0);	\coordinate (B) at (6,0);|\newline
%    \verb|      \coordinate (C) at (0,3);|\newline
%    \verb|      \draw (A) -- (B) -- node[above right]{hypoténuse} (C) -- cycle;|\newline
%    \verb|      \draw (0,0) rectangle (0.3,0.3);|\newline
%    \verb|    \end{tikzpicture}};|\newline
%    \verb|  \node [annotation,above right=10pt] at (hypotenuse.north east){|\newline
%    \verb|    C'est le plus grand côté du triangle.\newline |\newline
%    \verb|    C'est le côté à l'opposé de l'angle droit.\newline|\newline
%    \verb|    Seuls les triangles rectangles ont une hypoténuse.};|\newline
%    \verb|  \node [annotation,left=5pt] at (theoreme.north west){|\newline
%    \verb|    Si un triangle est rectangle, alors le carré de la longueur de l'hypoténuse |\newline
%    \verb|    est égal à la somme des carrés des longueurs des deux autres côtés.};|\newline
%    \verb|  \node [annotation,below left=5pt] at (theoreme.south west){|\newline
%    \verb|    Si, dans un triangle, le carré de la longueur du plus grand côté|\newline
%    \verb|    n'est pas égal à la somme des carrés des longueurs de deux autres côtés,|\newline
%    \verb|    alors le triangle n'est pas rectangle.\newline|\newline
%    \verb|    \textit{(Cet énoncé s'appelle la contraposée du théorème)}};|\newline
%    \verb|  \node [annotation,below right=5pt] at (reciproque.south east){|\newline
%    \verb|    Si, dans un triangle, le carré de la longueur du plus grand côté est égal|\newline
%    \verb|    à la somme des carrés des longueurs des deux autres côtés,|\newline
%    \verb|    alors le triangle est rectangle.};|\newline
%    \verb|  \node [annotation,below right=5pt] at (egalite.east){|\newline
%    \verb|    On a $BC^2=AB^2+AC^2$, c'est à dire \newline|\newline
%    \verb|    $BC=\sqrt{AB^2+AC^2}$ et $AB=\sqrt{BC^2-AC^2}$ et $AC=\sqrt{BC^2-AB^2}$};|\newline
%    \verb|  \node [annotation,above right=5pt] at (egalite.east){|\newline
%    \verb|    \begin{tikzpicture}[sharp corners]|\newline
%    \verb|      \coordinate (A) at (0,0);	\coordinate (B) at (2,0);|\newline
%    \verb|      \coordinate (C) at (0,1);|\newline
%    \verb|      \draw (A) -- (B) -- (C) -- cycle;|\newline
%    \verb|      \node at (-0.5,-0.5) {$A$};	\node at (2,-0.5) {$B$};|\newline
%    \verb|      \node at (-0.5,1) {$C$};|\newline
%    \verb|      \draw (0,0) rectangle (0.2,0.2);|\newline
%    \verb|    \end{tikzpicture}|\newline
%    \verb|  };|\newline
%    \verb| \end{tikzpicture} \end{center}|\newline
%    \verb|\end{seqCarteMentale}|\newline
%    donne :
%    \begin{seqCarteMentale}%
%     \begin{center}%
%      \begin{tikzpicture}[small mindmap,concept color=\seqCouleurVioletVif,
%           level 1 concept/.append style={concept color=\seqCouleurVioletModere},
%           every annotation/.append style={fill=\seqCouleurVioletPale}]
%    	 \node [concept] {Théorème de Pythagore} [counterclockwise from=0]
%           child {node[concept] (egalite) {\'Egalité de Pythagore} }%
%           child {node[concept] (hypotenuse) {Hypoténuse d'un triangle rectangle} }%
%           child {node[concept] (racine) {Racine carrée} }%
%           child {node[concept] (theoreme) {Le théorème} }%
%           child {node[concept] (reciproque) {La réciproque du théorème} };
%        \node [annotation,above left=5pt] at (racine.north west){%
%           La racine carrée du nombre positif $a$ est le nombre positif $b$ tel que $b^2=a$.};
%        \node [annotation,left=5pt] at (racine.west){%
%           Le symbole de la  racine carrée est $\sqrt{\phantom{ab}}$.};
%        \node [annotation,above left=5pt] at (hypotenuse.north){%
%           \begin{tikzpicture}[scale=0.4, sharp corners]
%              \coordinate (A) at (0,0);	\coordinate (B) at (6,0);	\coordinate (C) at (0,3);
%              \draw (A) -- (B) -- node[above right]{hypoténuse} (C) -- cycle;
%              \draw (0,0) rectangle (0.3,0.3);
%           \end{tikzpicture}};
%        \node [annotation,above right=10pt] at (hypotenuse.north east){%
%           C'est le plus grand côté du triangle.\newline C'est le côté à l'opposé de l'angle droit.\newline
%           Seuls les triangles rectangles ont une hypoténuse.};
%        \node [annotation,left=5pt] at (theoreme.north west){%
%           Si un triangle est rectangle, alors le carré de la longueur de l'hypoténuse est égal 
%           à la somme des carrés des longueurs des deux autres côtés.};
%        \node [annotation,below left=5pt] at (theoreme.south west){%
%           Si, dans un triangle, le carré de la longueur du plus grand côté n'est pas égal 
%           à la somme des carrés des longueurs des deux autres côtés, alors le triangle n'est pas rectangle.\newline
%           \textit{(Cet énoncé s'appelle la contraposée du théorème)}};
%        \node [annotation,below right=5pt] at (reciproque.south east){%
%           Si, dans un triangle, le carré de la longueur du plus grand côté est égal à la somme des carrés
%           des longueurs de deux autres côtés, alors le triangle est rectangle.};
%        \node [annotation,below right=5pt] at (egalite.east){%
%           On a $BC^2=AB^2+AC^2$, c'est à dire \newline
%           $BC=\sqrt{AB^2+AC^2}$ et $AB=\sqrt{BC^2-AC^2}$ et $AC=\sqrt{BC^2-AB^2}$};
%        \node [annotation,above right=5pt] at (egalite.east){%
%           \begin{tikzpicture}[sharp corners]
%              \coordinate (A) at (0,0);	\coordinate (B) at (2,0);	\coordinate (C) at (0,1);
%              \draw (A) -- (B) -- (C) -- cycle;
%              \node at (-0.5,-0.5) {$A$};	\node at (2,-0.5) {$B$};	\node at (-0.5,1) {$C$};
%              \draw (0,0) rectangle (0.2,0.2);
%           \end{tikzpicture}%
%       };
%     \end{tikzpicture} \end{center}%
%    \end{seqCarteMentale}%
%
%  \subsection{Méthodes clés}%
%   \DescribeEnv{methode}%
%    L'environnement |methode|\marg{titre}\marg{enonce} construit une boîte comportant un titre et contenant deux parties:
%    l'énoncé des exercices tout d'abord, et leurs corrigés ensuite.\newline
%    L'argument obligatoire \meta{titre} est le titre de la méthode. \newline
%    L'argument obligatoire \meta{enonce} est l'enoncé des exercices.\newline
%    Les méthodes sont numérotées automatiquement.
%    \medskip \newline
%   Par exemple:\newline
%   \verb|  \begin{seqMethode}{Calculer la longueur de l'hypoténuse d'un triangle rectangle}{|\newline
%   \verb|    Le triangle $ABC$ est un triangle rectangle en $A$ tel que $AB=12$ cm|\newline
%   \verb|    et $AC=5$ cm. \newline Calculer la longueur de $[BC]$.|\newline
%   \verb|   }|\newline
%   \verb|   \begin{multicols}{2}|\newline
%   \verb|    Je commence par faire une figure.\newline|\newline
%   \verb|    Le triangle $ABC$ étant rectangle en $A$, le théorème de Pythagore permet |\newline
%   \verb|    d'écrire l'égalité de Pythagore: \newline $AB^2+AC^2=BC^2$\newline|\newline
%   \verb|    Je calcule $AB^2=12^2=144$ et $AC^2=5^2=25$. |\newline
%   \verb|    Ainsi, $BC^2=144+25=169$.\newline|\newline
%   \verb|    Je calcule ensuite la racine carrée de 169 avec la calculatrice: | \newline
%   \verb|    $\sqrt{169}=13$.\newline|\newline
%   \verb|    Finalement, $BC=13$ cm.|\newline
%   \verb|    \begin{center} \begin{tikzpicture}[scale=0.35]|\newline
%   \verb|     \coordinate (A) at (0,0); \coordinate (B) at (12,0); |\newline
%   \verb|     \coordinate (C) at (0,5);|\newline
%   \verb|     \draw (A) node[below left]{$A$} -- node[below]{$12$ cm}|\newline
%   \verb|       (B) node[below right]{$B$} -- node[above right]{longueur à calculer}|\newline
%   \verb|       (C) node[above]{$C$} -- node[left]{$5$ cm} cycle;|\newline
%   \verb|     \draw (0,0) rectangle (0.4,0.4);|\newline
%   \verb|    \end{tikzpicture} \end{center}|\newline
%   \verb|   \end{multicols}|\newline
%   \verb|  \end{seqMethode}|\newline
%   donne :
%   \begin{seqMethode}{Calculer la longueur de l'hypoténuse d'un triangle rectangle}{%
%     Le triangle $ABC$ est un triangle rectangle en $A$ tel que $AB=12$ cm et $AC=5$ cm.\newline Calculer la longueur de $[BC]$.
%    }%
%    \begin{multicols}{2}%
%     Je commence par faire une figure.\newline
%     Le triangle $ABC$ étant rectangle en $A$, le théorème de Pythagore permet d'écrire l'égalité de Pythagore:\newline
%     $AB^2+AC^2=BC^2$\newline
%     Je calcule $AB^2=12^2=144$ et $AC^2=5^2=25$. Ainsi, $BC^2=144+25=169$.\newline
%     Je calcule ensuite la racine carrée de 169 avec la calculatrice: $\sqrt{169}=13$.\newline
%     Finalement, $BC=13$ cm.
%     \begin{center} \begin{tikzpicture}[scale=0.35]
%      \coordinate (A) at (0,0);	\coordinate (B) at (12,0);	\coordinate (C) at (0,5);
%      \draw (A) node[below left]{$A$} -- node[below]{$12$ cm} (B) node[below right]{$B$} -- 
%       node[above right]{longueur à calculer} (C) node[above]{$C$} -- node[left]{$5$ cm} cycle;
%      \draw (0,0) rectangle (0.4,0.4);
%     \end{tikzpicture} \end{center}%
%    \end{multicols}%
%   \end{seqMethode}%
%
%  \subsection{Séries d'exercices et exercices}%
%   \DescribeEnv{serieExos}%
%    L'environnement |serieExos|\oarg{titre} construit une boîte comportant un titre, dans laquelle seront regroupés les exercices correspondants. Il saute une page quand il se termine. \newline
%    L'argument optionnel \meta{titre} est le titre de la série d'exercices.
%    \medskip \newline
%   \DescribeEnv{exercice}%
%    L'environnement |exercice|\oarg{titre} numérote automatiquement l'exercice dans la série à laquelle il appartient.\newline
%    L'argument optionnel \meta{titre} est le titre de l'exercice, vide par défaut.
%    \medskip \newline
%   \DescribeMacro{\seqCorrige}%
%    La macro |\seqCorrige|\marg{corrige} permet de donner le corrigé de l'exercice. Elle se place à l'intérieur de l'environnement |\exercice|.\newline
%    L'argument obligatoire \meta{corrige} est le texte du corrigé.\newline
%    Tous les corrigés sont automatiquement regroupés en fin de document.
%    \medskip \newline
%   Par exemple:\newline
%   \verb|  \begin{seqSerieExos}{fondamentale}|\newline
%   \verb|   \begin{seqExercice}|\newline
%   \verb|    Calculer les nombres suivants.|\newline
%   \verb|    \begin{seqColEnum}[3]|\newline
%   \verb|     \item $3^2$ \item $10^2$ \item $9^2$ \item $4^2+5^2$|\newline
%   \verb|     \item $7^2+3^2$ \item $9^2+6^2$|\newline
%   \verb|    \end{seqColEnum}|\newline
%   \verb|    \seqCorrige{|\newline
%   \verb|     Je calcule les résultats des expressions :|\newline
%   \verb|     \begin{seqColEnum}[3]|\newline
%   \verb|      \item $3^2 = 3 \times 3 = 9$ \item $10^2 = 10 \times 10 = 100$|\newline
%   \verb|      \item $9^2 = 9 \times 9 = 81$ \item $4^2 + 5^2 = 16 + 25 = 41$|\newline
%   \verb|      \item $7^2 + 3^2 = 49 + 9 = 58$ \item $9^2 + 6^2 = 81 + 36 = 117$|\newline
%   \verb|     \end{seqColEnum}|\newline
%   \verb|    }|\newline
%   \verb|   \end{seqExercice}|\newline
%   \verb|   \begin{seqExercice}|\newline
%   \verb|    Trouver tous les triangles rectangles dans la figure ci-dessous.\newline|\newline
%   \verb|    Pour chacun, nommer son hypoténuse.|\newline
%   \verb|    \begin{center} \begin{tikzpicture}[scale=0.5]|\newline
%   \verb|     \tkzInit[ymin=-1,ymax=7,xmin=-1,xmax=10] \tkzClip|\newline
%   \verb|     \tkzDefPoint(0,0){C} \tkzDefPoint(6,0){B} \tkzDefPoint(0,3){A}|\newline
%   \verb|     \tkzDefLine[perpendicular=through C](A,B) \tkzGetPoint{m}|\newline
%   \verb|     \tkzInterLL(C,m)(A,B) \tkzGetPoint{F}|\newline
%   \verb|     \tkzDefPointBy[translation= from C to F](A) \tkzGetPoint{D}|\newline
%   \verb|     \tkzDefPointBy[rotation=center B angle -90](D) \tkzGetPoint{n}|\newline
%   \verb|     \tkzDefPointBy[homothety=center B ratio .4](n) \tkzGetPoint{E}|\newline
%   \verb|     \tkzDrawPolygon(A,B,C) \tkzDrawPolygon(A,B,D)|\newline
%   \verb|     \tkzDrawPolygon(A,F,C) \tkzDrawPolygon(B,D,E)|\newline
%   \verb|     \tkzMarkRightAngle(B,C,A) \tkzMarkRightAngle(B,A,D)|\newline
%   \verb|     \tkzMarkRightAngle(A,F,C) \tkzMarkRightAngle(E,B,D)|\newline
%   \verb|     \tkzLabelPoints[above](D,F) \tkzLabelPoints[below](B,C)|\newline
%   \verb|     \tkzLabelPoints[left](A) \tkzLabelPoints[right](E)|\newline
%   \verb|    \end{tikzpicture} \end{center}|\newline
%   \verb|    \seqCorrige{|\newline
%   \verb|     Le triangle $ABC$ a pour hypoténuse $AB$. \quad|\newline
%   \verb|     Le triangle $ABD$ a pour hypoténuse $DB$. \newline|\newline
%   \verb|     Le triangle $BDE$ a pour hypoténuse $DE$. \quad|\newline
%   \verb|     Le triangle $AFC$ a pour hypoténuse $AC$. \newline|\newline
%   \verb|     Le triangle $BFC$ a pour hypoténuse $BC$.|\newline
%   \verb|    }|\newline
%   \verb|   \end{seqExercice}|\newline
%   \verb|  \end{seqSerieExos}|\newline
%   donne : 
%   \newpage
%   \begin{seqSerieExos}{fondamentale}%
%    \begin{seqExercice}[nom=carrés]
%     Calculer les nombres suivants.
%     \begin{seqColEnum}[nbCols=3]
%      \item $3^2$	\item $10^2$	\item $9^2$	\item $4^2+5^2$
%      \item $7^2+3^2$	\item $9^2+6^2$
%     \end{seqColEnum}%
%     \seqCorrige{%
%      Je calcule les résultats des expressions :
%      \begin{seqColEnum}[nbCols=3]
%       \item $3^2 = 3 \times 3 = 9$
%       \item $10^2 = 10 \times 10 = 100$
%       \item $9^2 = 9 \times 9 = 81$
%       \item $4^2 + 5^2 = 16 + 25 = 41$
%       \item $7^2 + 3^2 = 49 + 9 = 58$
%       \item $9^2 + 6^2 = 81 + 36 = 117$
%      \end{seqColEnum}%
%     }%
%    \end{seqExercice}%
%    \begin{seqExercice}%
%     Trouver tous les triangles rectangles dans la figure ci-dessous. \newline Pour chacun, nommer son hypoténuse.
%     \begin{center} \begin{tikzpicture}[scale=0.5]
%      \tkzInit[ymin=-1,ymax=7,xmin=-1,xmax=10] \tkzClip
%      \tkzDefPoint(0,0){C} \tkzDefPoint(6,0){B} \tkzDefPoint(0,3){A}%
%      \tkzDefLine[perpendicular=through C](A,B) \tkzGetPoint{m}%
%      \tkzInterLL(C,m)(A,B) \tkzGetPoint{F}%
%      \tkzDefPointBy[translation= from C to F](A) \tkzGetPoint{D}%
%      \tkzDefPointBy[rotation=center B angle -90](D) \tkzGetPoint{n}%
%      \tkzDefPointBy[homothety=center B ratio .4](n) \tkzGetPoint{E}%
%      \tkzDrawPolygon(A,B,C)	\tkzDrawPolygon(A,B,D)	\tkzDrawPolygon(A,F,C)	\tkzDrawPolygon(B,D,E)
%      \tkzMarkRightAngle(B,C,A)	\tkzMarkRightAngle(B,A,D)	\tkzMarkRightAngle(A,F,C)	\tkzMarkRightAngle(E,B,D)
%      \tkzLabelPoints[above](D,F)	\tkzLabelPoints[below](B,C)	\tkzLabelPoints[left](A)	\tkzLabelPoints[right](E)
%     \end{tikzpicture} \end{center}%
%    \seqCorrige{%
%     Le triangle $ABC$ a pour hypoténuse $AB$. \quad Le triangle $ABD$ a pour hypoténuse $DB$. \newline
%     Le triangle $BDE$ a pour hypoténuse $DE$. \quad Le triangle $AFC$ a pour hypoténuse $AC$. \newline
%     Le triangle $BFC$ a pour hypoténuse $BC$.
%    }%
%   \end{seqExercice}%
%  \end{seqSerieExos}%
%
%  \subsection{Affichage des corrigés}%
%   \DescribeMacro{\seqAfficheCorriges}%
%    La macro |\seqAfficheCorriges| permet d'afficher les corrigés de tous les exercices, y compris les révisions.
%    Elle se place à la fin du document.
%    \medskip \newline
%    Exemple avec les exercices de ce document:
%    \seqAfficheCorriges
%
% \newpage
% \section{Les différents éléments de l'évaluation bilan} \label{subsec:evalOutils} \label{elementsEval}%
%  \subsection{Titre de l'évaluation bilan}%
%   \DescribeMacro{\titreEval}%
%    La macro |\titreEval|\oarg{options} met en forme le titre de l'évaluation.\newline
%    Elle définit les entêtes et pieds de page.\newline
%    Elle crée les compteurs pour numéroter automatiquement les exercices.\newline
%    L'option \texttt{numSeq}=\meta{texte} est le numéro de la séquence dans la progression. \newline
%    L'option \texttt{titreSeq}=\meta{texte} est le nom de la séquence. \newline
%    L'option \texttt{niveauDiff}=\meta{texte} est le niveau de difficulté de cette évaluation. Elle peut prendre les valeurs \texttt{fondamental}, \texttt{avancé} ou \texttt{exploration}. \newline
%    L'option \texttt{etab}=\meta{texte} est le nom de l'établissement. \newline
%    L'option \texttt{classe}=\meta{texte} est le nom de la classe. \newline
%    Ces cinq options sont obligatoires. \newline
%    L'option facultative \texttt{theme}=\meta{texte} définit le thème de couleurs utilisé.
%     Les différentes valeurs possibles sont \textit{nombres}, \textit{donnees}, \textit{geometrie} et \textit{gris} (valeur par défaut). 
%     Voir \ref{subsec:themes} pour leur définition.
%   \medskip \newline
%   Par exemple : \newline
%   \verb|\titreEval[numSeq=01,titreSeq=Proportionnalité,niveauDiff=fondamental]| \newline
%   \verb|           etab=collège Exemple,classe=4\ieme{}A,theme=nombres]| \newline
%   donne
%	\par\vspace{20pt}%
%	\titreEval[numSeq=01,titreSeq=Proportionnalité,niveauDiff=fondamental,
%              etab=collège Exemple,classe=4\ieme{}A,theme=nombres]
%
%  \subsection{Barème de l'évaluation bilan}%
%   \DescribeEnv{bareme}%
%    L'environnement |bareme| construit une boîte comportant un titre et une courte description, 
%    dans laquelle seront listés les prérequis de la séquence en utilisant la macro |\baremeItem|.\newline
%    L'argument optionnel \meta{description} permet de modifier la description par défaut.
%   \medskip \newline
%   \DescribeMacro{\baremeItem}%
%    La macro |\baremeItem|\marg{nom}\marg{points} met en forme une ligne du barème.\newline
%    L'argument obligatoire \meta{nom} est le nom de l'exercice. \newline
%    L'argument obligatoire \meta{points} est le nombre de points de l'exercice. 
%    \medskip \newline
%   Par exemple : \newline
%   \verb|\begin{bareme}| \newline
%   \verb|	\baremeItem{QCM}{5 points}| \newline
%   \verb|	\baremeItem{Exercice II}{6 points}| \newline
%   \verb|	\baremeItem{Présentation de la copie et usage de la langue française}{1 point}| \newline
%   \verb|\end{bareme}| \newline
%   donne
%	\par\vspace{20pt}%
%   \begin{bareme}%
%   	\baremeItem{QCM}{5 points}%
%   	\baremeItem{Exercice II}{6 points}%
%   	\baremeItem{Présentation de la copie et usage de la langue française}{1 point}%
%   \end{bareme}%
%
%  \subsection{Exercices de l'évaluation bilan}%
%   \subsubsection{Corrigé d'un exercice} \label{subsec:corrExoEval}%
%    \DescribeMacro{\corrigeExoEval}%
%    La macro |\corrigeExoEval|\marg{corrige} permet de donner le corrigé de l'exercice. \newline
%    Elle se place à l'intérieur de l'environnement |qcm| ou |exerciceEval|.\newline
%    L'argument obligatoire \meta{corrige} est le texte du corrigé.\newline
%    Tous les corrigés sont automatiquement regroupés en fin d'évaluation. \newline
%    Des exemples d'utilisation sont donnés dans le paragraphes suivants.
%
%   \subsubsection{QCM}%
%   Les QCM sont mis en forme de manière particulière par rapport aux autres exercices, par l'utilisation des macros et environnements suivants.
%   \medskip \newline
%   \DescribeEnv{qcm}%
%    L'environnement |qcm|\oarg{options} met en forme le titre du qcm ainsi que les instructions pour y répondre.\newline
%    Il débute également un tableau dans lequel on placera les différentes questions composant le qcm en utilisant la macro |\qcmItem|.\newline
%    Les réponses possibles sont automatiquement nommées \texttt{A}, \texttt{B}, etc. \newline
%    Le corrigé du QCM est donné en utilisant la macro |\corrigeExoEval| décrite ci-dessus. \newline
%    L'option \texttt{ptOK}=\meta{texte} est le nombre de point(s) que rapporte une bonne réponse. \newline
%    L'option \texttt{ptKO}=\meta{texte} est le nombre de point(s) que rapporte une réponse partielle ou qu'enlève une mauvaise réponse. \newline
%    L'option \texttt{nbReps}=\meta{nombre} est le nombre de réponses pour chaque question. Il vaut 3 par défaut.\newline
%    L'option \texttt{largeurQuestion}=\meta{nombre} est la pondération de la largeur de la colonne pour la question. Elle vaut 3 par défaut. \newline
%    L'option \texttt{largeurRep}=\meta{nombre} est la pondération de la largeur de chaque colonne de réponse. Elle vaut 1 par défaut. \newline
%    Les options \texttt{ptOK} et \texttt{ptKO} sont obligatoires.
%   \medskip \newline
%   \DescribeMacro{\qcmItem}%
%    La macro |\qcmItem|\marg{enonce} met en forme une ligne du QCM.\newline
%    L'argument obligatoire \meta{enonce} est l'énoncé complet : il comprend la question et toutes les réponses possibles.
%    Chaque élément doit être séparé du suivant par un \&.
%    \medskip \newline
%   Par exemple : \newline
%   \verb|\begin{qcm}[ptOK=un point,ptPartiel=un demi-point,ptKO=un quart de point]| \newline
%   \verb| \qcmItem{| \newline
%   \verb|  Quels tableaux représentent une situation de proportionnalité ? &| \newline
%   \verb|  \begin{tblr}{X[r] *{2}{X[r]}}| \newline
%   \verb|   \hline 0 & 5 & 7 \\ \hline 6 & 30 & 42 \\ \hline| \newline
%   \verb|  \end{tblr} &| \newline
%   \verb|  \begin{tblr}{X[r] *{2}{X[r]}}| \newline
%   \verb|   \hline 6 & 8 & 14 \\ \hline 9 & 12 & 21 \\ \hline| \newline
%   \verb|  \end{tblr} &| \newline
%   \verb|  \begin{tblr}{X[r] *{2}{X[r]}}| \newline
%   \verb|   \hline 5 & 6 & 7 \\ \hline 8 & 9,6 & 11,2 \\ \hline| \newline
%   \verb|  \end{tblr}| \newline
%   \verb| }| \newline
%   \verb| \corrigeExoEval{| \newline
%   \verb|  \begin{seqColEnum}[nbCols=1]| \newline
%   \verb|   \item Les réponses \textbf{B} et \textbf{C} sont correctes. \newline| \newline
%   \verb|    Pour \textbf{A} : | \newline
%   \verb|     sur la première colonne, 0, et pas 6, devrait correspondre à 0. \newline| \newline
%   \verb|    Pour \textbf{B} : | \newline
%   \verb|     on passe de la première à la deuxième ligne en multipliant par | \newline
%   \verb|     1,5 pour toutes les colonnes. \newline| \newline
%   \verb|    Pour \textbf{C} : | \newline
%   \verb|     on passe de la première à la deuxième ligne en multipliant par | \newline
%   \verb|     1,6 pour toutes les colonnes. \newline| \newline
%   \verb|  \end{seqColEnum}| \newline
%   \verb| }| \newline
%   \verb|\end{qcm}| \newline
%   donne :
%   \begin{qcm}[ptOK=un point,ptPartiel=un demi-point,ptKO=un quart de point,nbReps=3,largeurQuestion=2,largeurRep=1]
%   	\qcmItem{%
%   		Quels tableaux représentent une situation de proportionnalité ? &
%   		\begin{tblr}{X[r] *{2}{X[r]}} \hline 0 & 5 & 7 \\ \hline 6 & 30 & 42 \\ \hline \end{tblr} &
%   		\begin{tblr}{X[r] *{2}{X[r]}} \hline 6 & 8 & 14 \\ \hline 9 & 12 & 21 \\ \hline \end{tblr} &
%   		\begin{tblr}{X[r] *{2}{X[r]}} \hline 5 & 6 & 7 \\ \hline 8 & 9,6 & 11,2 \\ \hline \end{tblr}%
%   	}%
%		\corrigeExoEval{%
%   		\begin{seqColEnum}[nbCols=1]
%   			\item Les réponses \textbf{B} et \textbf{C} sont correctes. \newline
%   				Pour \textbf{A} : sur la première colonne, 0, et pas 6, devrait correspondre à 0. \newline
%   				Pour \textbf{B} : on passe de la première à la deuxième ligne en multipliant par 1,5 pour toutes les colonnes. \newline
%   				Pour \textbf{C} : on passe de la première à la deuxième ligne en multipliant par 1,6 pour toutes les colonnes. \newline
%   		\end{seqColEnum}%
%   	}%
%   \end{qcm}%
%
%   \subsubsection{Autres exercices}%
%    Les autres exercices de l'évaluation sont mis en forme avec l'environnement suivant :
%    \medskip \newline
%    \DescribeEnv{exerciceEval}%
%     L'environnement |exerciceEval|\oarg{titre} numérote automatiquement l'exercice. \newline
%     Il insere une boîte de titre au début de l'exercice. \newline
%     L'argument \meta{titre} est facultatif.
%     Le corrigé du QCM est donné en utilisant la macro |\corrigeExoEval| décrite au \ref{subsec:corrExoEval}.
%     \medskip \newline
%    Par exemple : \newline
%    \verb|\begin{exerciceEval}[nom=avec les pieds]| \newline
%    \verb| Le pied (en anglais \textit{foot} noté ft) est une unité de longueur | \newline
%    \verb| anglo-saxonne : 5000\,ft correspondent à 1524\,m.| \newline
%    \verb| \begin{seqColEnum}| \newline
%    \verb|  \item Un ULM vole à 800\,m d'altitude. Convertir cette altitude en pieds. | \newline
%    \verb|     Donner une valeur approchée à l'unité près.| \newline
%    \verb|   \begin{center} | \newline
%    \verb|    \includegraphics[scale=0.2]{Images/EvalFondamentalExo2Enonce.jpg}| \newline
%    \verb|   \end{center}| \newline
%    \verb|  \item Pour les appareils qui effectuent des vols \og à vue \fg{}, la limite | \newline
%    \verb|    de survol de certaines villes est fixée à 3300\,ft.| \newline
%    \verb|    Convertir cette altitude en mètres.| \newline
%    \verb|  \item Pour effectuer les calculs plus facilement, Noah utilise 30\,cm comme | \newline
%    \verb|    correspondance pour 1\,ft.| \newline
%    \verb|    Pour une altitude de 200\,ft, quelle est la différence, en cm, | \newline
%    \verb|    entre la mesure de Noah et la valeur exacte ?| \newline
%    \verb| \end{seqColEnum}| \newline
%    \verb| \corrigeExoEval{| \newline
%    \verb|  \begin{seqColEnum}| \newline
%    \verb|   \item Pour convertir la valeur, j'utilise un tableau de proportionnalité : | \newline
%    \verb|    \begin{center} \begin{tblr}{X[2,r] *{2}{X[r]}}| \newline
%    \verb|     \hline En mètres & 1524 & 800 \\ \hline En pieds & 5000 & $x$ \\ \hline| \newline
%    \verb|    \end{tblr} \end{center}| \newline
%    \verb|    Je calcule en utilisant le produit en croix : | \newline
%    \verb|    $x=5000 \times 800 \div 1524 \approx 2625$. \newline| \newline
%    \verb|    800\;m correspond environ à 2625\;ft.| \newline
%    \verb|   \item Pour convertir la valeur, j'utilise un tableau de proportionnalité : | \newline
%    \verb|    \begin{center} \begin{tblr}{X[2,r] *{2}{X[r]}}| \newline
%    \verb|     \hline En mètres & 1524 & $x$ \\ \hline En pieds & 5000 & 3300 \\ \hline| \newline
%    \verb|    \end{tblr} \end{center}| \newline
%    \verb|    Je calcule en utilisant le produit en croix : | \newline
%    \verb|    $x=3300 \times 1524 \div 5000 = 1005,84$. \newline| \newline
%    \verb|    3300\;ft correspond à 1005,84\;m.| \newline
%    \verb|   \item Pour convertir la valeur réelle, j'utilise un tableau de | \newline
%    \verb|    proportionnalité : | \newline
%    \verb|    \begin{center} \begin{tblr}{X[2,r] *{2}{X[r]}}| \newline
%    \verb|     \hline En mètres & 1524 & $x$ \\ \hline En pieds & 5000 & 200 \\ \hline| \newline
%    \verb|    \end{tblr} \end{center}| \newline
%    \verb|    Je calcule en utilisant le produit en croix : | \newline
%    \verb|    $x=200 \times 1524 \div 5000 = 60,96$. \newline| \newline
%    \verb|    200\;ft correspond à 60,96\;m soit 6096\;cm. \newline| \newline
%    \verb|    J'utilise maintenant la méthode de Noah : $30 \times 200 = 6000$.| \newline
%    \verb|    Pour lui, 200\;ft correspond à 6000\;cm. \newline| \newline
%    \verb|    Il y a 96\;cm d'écart entre les deux valeurs.| \newline
%    \verb|  \end{seqColEnum}| \newline
%    \verb| }| \newline
%    \verb|\end{exerciceEval}| \newline
%    donne
%    \begin{exerciceEval}[nom=avec les pieds]
%    	Le pied (en anglais \textit{foot} noté ft) est une unité de longueur anglo-saxonne : 5000\,ft correspondent à 1524\,m.
%    		\begin{seqColEnum}%
%    			\item Un ULM vole à 800\,m d'altitude. Convertir cette altitude en pieds. Donner une valeur approchée à l'unité près.
%    			\item Pour les appareils qui effectuent des vols \og à vue \fg{}, la limite de survol de certaines villes est fixée à 3300\,ft.
%    					Convertir cette altitude en mètres.
%    			\item Pour effectuer les calculs plus facilement, Noah utilise 30\,cm comme correspondance pour 1\,ft.
%    					Pour une altitude de 200\,ft, quelle est la différence, en cm, entre la mesure de Noah et la valeur exacte ?
%    		\end{seqColEnum}%
%    	\corrigeExoEval{%
%        \begin{seqColEnum}%
%         \item Pour convertir la valeur, j'utilise un tableau de proportionnalité : 
%          \begin{center} \begin{tblr}{X[2,r] *{2}{X[r]}}%
%           \hline En mètres & 1524 & 800 \\ \hline En pieds & 5000 & $x$ \\ \hline
%          \end{tblr} \end{center}%
%          Je calcule en utilisant le produit en croix : $x=5000 \times 800 \div 1524 \approx 2625$. \newline
%          800\;m correspond environ à 2625\;ft.
%         \item Pour convertir la valeur, j'utilise un tableau de proportionnalité : 
%          \begin{center} \begin{tblr}{X[2,r] *{2}{X[r]}}%
%           \hline En mètres & 1524 & $x$ \\ \hline En pieds & 5000 & 3300 \\ \hline
%          \end{tblr} \end{center}%
%          Je calcule en utilisant le produit en croix : $x=3300 \times 1524 \div 5000 = 1005,84$. \newline
%          3300\;ft correspond à 1005,84\;m.
%         \item Pour convertir la valeur réelle, j'utilise un tableau de proportionnalité : 
%          \begin{center} \begin{tblr}{X[2,r] *{2}{X[r]}}%
%           \hline En mètres & 1524 & $x$ \\ \hline En pieds & 5000 & 200 \\ \hline
%          \end{tblr} \end{center}%
%          Je calcule en utilisant le produit en croix : $x=200 \times 1524 \div 5000 = 60,96$. \newline
%          200\;ft correspond à 60,96\;m soit 6096\;cm. \newline
%          J'utilise maintenant la méthode de Noah : $30 \times 200 = 6000$.
%          Pour lui, 200\;ft correspond à 6000\;cm. \newline
%          Il y a 96\;cm d'écart entre les deux valeurs.
%        \end{seqColEnum}%
%    	}%
%    \end{exerciceEval}%
%
%  \subsection{Affichage du corrigé de l'évaluation}%
%   \DescribeMacro{\afficheCorrigesEval}%
%    La macro |\afficheCorrigesEval| permet d'afficher les corrigés de tous les exercices, y compris les QCM, sur une nouvelle page impaire.
%    Elle se place à la fin de l'évaluation.
%    \medskip \newline
%    Exemple avec les exercices de cette évaluation sur la prochaine page impaire :
%    \afficheCorrigesEval

% \newpage
% \section{Chargement des données pédagogiques} \label{sec:donnees}%
%  Le paquet \texttt{seqenseigne} peut exploiter des bases de données pédagogiques
%  (objectifs, connaissances, séquences) stockées dans des fichiers \texttt{.dbtex}%
%  précompilés avec l'outil \texttt{datatooltk-gui}.
%  Cette fonctionnalité est optionnelle : les macros structurelles du livret et de
%  l'évaluation fonctionnent sans elle, mais les macros \texttt{Get*} et les tableaux
%  automatiques (tableau d'objectifs, tableau de prérequis) l'exigent.
%
%  \subsection{Organisation des fichiers de données}%
%   Les bases sont organisées par cycle. Les fichiers attendus sont :
%   \begin{seqColItem}[nbCols=1]
%    \item \texttt{param\_niveaux.dbtex} : correspondance codes niveaux / cycles
%    \item \texttt{C03\_themes.dbtex}, \texttt{C03\_sequences.dbtex},
%          \texttt{C03\_connaissances.dbtex}, \texttt{C03\_objectifs.dbtex}%
%    \item \texttt{C04\_themes.dbtex}, \texttt{C04\_sequences.dbtex},
%          \texttt{C04\_attendus\_fin\_C03.dbtex},
%          \texttt{C04\_connaissances.dbtex}, \texttt{C04\_objectifs.dbtex}%
%   \end{seqColItem}%
%   Par convention ces fichiers sont trois niveaux au-dessus du document compilé
%   (\texttt{../../../}). Ce chemin est modifiable avec |\seqSetDataPath|.
%
%  \subsection{Chargement principal}%
%   \DescribeMacro{\seqLoadData}%
%    |\seqLoadData| charge l'ensemble des bases cycles 3 et 4. Elle doit être appelée
%    \emph{une seule fois}, avant tout appel aux macros \texttt{Get*}.
%   \medskip \newline
%   \DescribeMacro{\seqSetDataPath}%
%    |\seqSetDataPath|\marg{chemin} modifie le répertoire de recherche des \texttt{.dbtex}.
%    Le \meta{chemin} doit se terminer par un \texttt{/}.\newline
%    Exemple pour un document au même niveau que les bases :\newline
%    \verb|\seqSetDataPath{./}| \newline
%    \verb|\seqLoadData|
%
%  \subsection{Macros d'accès}%
%   \DescribeMacro{\seqNiveauGetNomCourt}%
%    |\seqNiveauGetNomCourt|\marg{codeNiveau} : nom court (ex.~\texttt{N11} $\to$ \og quatrième \fg{}).
%   \medskip \newline
%   \DescribeMacro{\seqSequenceGetNom}%
%    |\seqSequenceGetNom|\marg{codeNiveau}\marg{codeSequence} : titre de la séquence.
%   \medskip \newline
%   \DescribeMacro{\seqSequenceGetNumero}%
%    |\seqSequenceGetNumero|\marg{codeNiveau}\marg{codeSequence} : numéro dans la progression.
%   \medskip \newline
%   \DescribeMacro{\seqObjectifGetNom}%
%    |\seqObjectifGetNom|\marg{codeNiveau}\marg{codeSequence}\marg{codeObjectif} : libellé d'un objectif.
%   \medskip \newline
%   \DescribeMacro{\seqConnaissanceGetNom}%
%    |\seqConnaissanceGetNom|\marg{codeNiveau}\marg{codeSequence}\marg{codeConnaissance} : libellé d'une connaissance.
%   \medskip \newline
%   \DescribeMacro{\seqSetColors}%
%    |\seqSetColors|\marg{codeNiveau}\marg{codeSequence} : règle le thème de couleurs
%    d'après le thème de la séquence dans les bases.
%
%  \subsection{Tableaux automatiques}%
%   \DescribeMacro{\seqTableauObjectifs}%
%    |\seqTableauObjectifs|\oarg{niveau=\meta{code},sequence=\meta{code}} : tableau
%    récapitulatif des critères de maîtrise (Très bon / Satisfaisant / À consolider).
%   \medskip \newline
%   \DescribeMacro{\seqTableauPrerequis}%
%    |\seqTableauPrerequis|\oarg{niveau=\meta{code},sequence=\meta{code}} : tableau
%    de prérequis adapté au niveau (attendus fin de cycle pour N10, objectifs des années
%    précédentes pour N08, N09, N11, N12).
%
% \newpage
% \section{Plan de travail automatique} \label{sec:plan}%
%  Le plan automatique exploite les bases pour générer objectifs et répartition des séances.
%
%  \subsection{Titre}%
%   \DescribeMacro{\seqBoiteTitrePlan}%
%    |\seqBoiteTitrePlan|\oarg{options} : boîte de titre utilisant |\seqCodeSequence|,
%    |\annee|, |\nomClasse| et |\nomCollege|.\newline
%    Options : \texttt{afficheSeq=oui|non}, \texttt{affichePlan=oui|non},
%    \texttt{titre=}\meta{texte} (obligatoire).
%
%  \subsection{Bloc objectif}%
%   \DescribeMacro{\seqPlanObjectif}%
%    |\seqPlanObjectif|\marg{niveau}\marg{sequence}\marg{objectif}\marg{connaissances}%
%    \marg{exosF}\marg{exosA}\marg{exosE} : bloc affichant le titre de l'objectif,
%    les notions et les numéros d'exercices des trois séries.\newline
%    Exemple :\newline
%    \verb|\seqPlanObjectif{N11}{S03}{02}| \newline
%    \verb|  {\seqConnaissanceGetNom{N11}{S03}{01}}{1,2,3}{4,5}{6}|
%
%  \subsection{Repères temporels}%
%   \DescribeMacro{\seqPlanReperesNiveaux}%
%    |\seqPlanReperesNiveaux|\oarg{options} : bloc repères temporels (dates, séances par
%    objectif) et bloc niveaux atteints, côte à côte.\newline
%    Options : \texttt{codeNiveau}, \texttt{codeSequence}, \texttt{dateDeb}, \texttt{dateFin},
%    \texttt{seancesTotal}, \texttt{seancesRevisions}, \texttt{seancesCours},
%    \texttt{seancesObj02} \ldots \texttt{seancesObj26}.
%
% \newpage
% \section{Suivi de classe} \label{sec:suivi}%
%
%  \subsection{Fichiers attendus}%
%   \begin{seqColItem}[nbCols=1]
%    \item \texttt{liste\_eleves.csv} : colonnes \texttt{Numero}, \texttt{Nom}, \texttt{Prenom}%
%    \item \texttt{periodes.csv} : colonne \texttt{Periode} (ex. \texttt{P1}, \texttt{P2}\ldots)
%    \item \texttt{P1-sequences.csv} \ldots : colonnes \texttt{Code}, \texttt{Nom}%
%    \item \texttt{P1bilan.csv} \ldots : colonne \texttt{Numero} + une colonne par séquence
%          contenant un code de maîtrise (\texttt{I}, \texttt{F}, \texttt{A}, \texttt{E}, \texttt{NE})
%    \item \texttt{S01Suivi.csv} \ldots : colonne \texttt{Numero} + une colonne par objectif
%   \end{seqColItem}%
%
%  \subsection{Chargement}%
%   \DescribeMacro{\seqLoadDataSuivi}%
%    |\seqLoadDataSuivi|\marg{codeNiveau} : charge les bases pédagogiques du cycle,
%    la liste des élèves, les périodes et tous les CSV de suivi. Remplace |\seqLoadData|
%    pour les documents de suivi.\newline Exemple : |\seqLoadDataSuivi{N11}|
%   \medskip \newline
%   \DescribeMacro{\seqTitreSuivi}%
%    |\seqTitreSuivi|\marg{classe}\marg{titre}\marg{annee} : configure en-têtes et pieds de page.
%
%  \subsection{Tableau annuel par élève}%
%   \DescribeMacro{\seqSuiviAnnuelParEleve}%
%    |\seqSuiviAnnuelParEleve|\oarg{niveau=\meta{code}} : pour chaque élève, tableau
%    paginé de ses niveaux de maîtrise par objectif sur toute l'année.\newline
%    Structure minimale :\newline
%    \verb|\documentclass[a4paper,11pt]{article}| \newline
%    \verb|\usepackage{seqenseigne}| \newline
%    \verb|\begin{document}| \newline
%    \verb|  \seqLoadDataSuivi{N11}| \newline
%    \verb|  \seqTitreSuivi{4\ieme{}A}{Suivi annuel}{2024/2025}| \newline
%    \verb|  \seqSuiviAnnuelParEleve[niveau=N11]| \newline
%    \verb|\end{document}|
%
%  \subsection{Tableau annuel par séquence}%
%   \DescribeMacro{\seqSuiviAnnuelParSequence}%
%    |\seqSuiviAnnuelParSequence|\oarg{niveau=\meta{code},sequence=\meta{code}} :
%    tableau d'une séquence, une colonne par élève, cellules colorées selon la maîtrise.
%   \medskip \newline
%   \DescribeMacro{\seqSuiviSequence}%
%    |\seqSuiviSequence|\marg{codeNiveau}\marg{codeSequence} : raccourci qui appelle
%    |\seqSuiviAnnuelParSequence| et passe à la page suivante.
%
%  \subsection{Tableau périodique par séquence}%
%   \DescribeMacro{\seqSuiviPeriodiqueParSequence}%
%    |\seqSuiviPeriodiqueParSequence|\oarg{niveau=\meta{code},periode=\meta{code}} :
%    bilan de période (une ligne par séquence, une colonne par élève) avec dates.
%
%  \subsection{Bilan périodique par élève}%
%   \DescribeMacro{\seqBilanPeriodiqueParEleve}%
%    |\seqBilanPeriodiqueParEleve|\oarg{options} : récapitulatif compact des maîtrises
%    par séquence pour chaque élève d'une période.\newline
%    Options : \texttt{periode}, \texttt{annee}, \texttt{niveau}, \texttt{nomClasse}.
%
%  \subsection{Codes de maîtrise}%
%   \begin{seqColItem}%
%    \item \texttt{I} : Insuffisant (rouge)
%    \item \texttt{F} : \`A consolider (jaune)
%    \item \texttt{A} : Satisfaisant (vert clair)
%    \item \texttt{E} : Très bon (vert foncé)
%    \item \texttt{NE} : Non évalué
%    \item \texttt{-} : Absent ou non concerné
%   \end{seqColItem}%

% Fin du module doc
%</package>
%
% \endinput
